% sections/09_governance.tex
\section{Governance and Versioning}\label{sec:governance}

\subsection{Custodian}\label{sec:gov:custodian}

The \ISI methodology is maintained by the \textbf{International Sovereignty
Institute\,---\,Research Unit}.  All methodology changes, snapshot
publications, and key-management decisions are the responsibility of the
custodian.

\subsection{Versioning Scheme}\label{sec:gov:versioning}

The \ISI follows a three-part version number: \texttt{major.minor.patch}.
The current version is \textbf{v\,1.0.0}.

\paragraph{Major version increment.}
A major version increment (e.g., 1.x.x~\(\to\)~2.0.0) is REQUIRED if any
of the following changes is made:

\begin{itemize}
  \item The aggregation rule (\cref{eq:composite}) changes.
  \item The number of axes changes (axes added or removed).
  \item Any classification threshold value (\cref{tab:thresholds}) changes.
  \item The rounding precision changes (currently 8~decimal places).
  \item The hashing algorithm changes (currently SHA-256).
  \item The signature algorithm changes (currently Ed25519).
  \item The canonical serialisation format
    (\cref{sec:det:serialisation}) changes.
  \item The hash-input field order or delimiter changes.
  \item The concentration formula (\cref{eq:hhi}) changes.
  \item The defence temporal window changes.
  \item The axis-weighting scheme changes.
\end{itemize}

A major-version change resets the minor and patch counters to~\(0\)
and requires a new methodology document.

\paragraph{Minor version increment.}
A minor version increment (e.g., 1.0.x~\(\to\)~1.1.0) is REQUIRED if
any of the following changes is made, provided no major-version trigger
applies:

\begin{itemize}
  \item Axis description text or metadata is revised without affecting
    the mathematical definition.
  \item The warning-code taxonomy (\cref{sec:appendix:warnings}) is expanded.
  \item The snapshot metadata schema (\cref{sec:snapshot:metadata}) is
    extended with new optional fields that do not enter the hash preimage.
  \item Country coverage is extended or reduced.
  \item A data-ingestion bug is corrected, producing different scores
    for one or more countries.
  \item A supplementary data source is added or replaced.
\end{itemize}

A minor-version change resets the patch counter to~\(0\).

\paragraph{Patch version increment.}
A patch version increment (e.g., 1.0.0~\(\to\)~1.0.1) is permitted
only if \emph{all} of the following conditions hold:

\begin{itemize}
  \item The change is limited to documentation corrections, typographic
    fixes, or non-normative commentary.
  \item No materialised output (score files, hash digests, signatures)
    is affected.
  \item The \texttt{methodology\_registry.json} semantics are unchanged.
  \item No snapshot is recomputed.
\end{itemize}

Patch versions do not alter the computation and therefore do not
require a new snapshot.

\subsection{Change Protocol}\label{sec:gov:change}

Proposed methodology changes follow a four-stage protocol:

\begin{enumerate}
  \item \textbf{Proposal.}  A written proposal specifies the change, its
    rationale, its impact on published scores (if any), and the affected
    version component (major or minor).
  \item \textbf{Impact assessment.}  The change is applied retroactively to
    the most recent two vintage years.  The resulting score differences are
    tabulated and reviewed.
  \item \textbf{Review.}  The proposal and impact assessment are reviewed by
    at least one independent methodologist not involved in the original
    proposal.
  \item \textbf{Publication.}  If approved, the change is incorporated into
    the next version release.  The methodology document, changelog, and
    version number are updated simultaneously.  The previous version remains
    archived and accessible.
\end{enumerate}

\subsection{Data-Revision Policy}\label{sec:gov:revision}

Statistical agencies routinely revise historical data.  When a source
revision affects the input data for a published vintage, the following rules
apply:

\begin{itemize}
  \item \textbf{No retroactive recalculation.}  Published snapshots are not
    recomputed.  The integrity chain binds each snapshot to the data as it
    existed at computation time.
  \item \textbf{Prospective incorporation.}  Revised data are incorporated
    into the next vintage computation.  The changelog records which source
    revisions were absorbed.
  \item \textbf{Erratum.}  If a source revision reveals a material error in
    a published snapshot (defined as a composite-score change exceeding
    \num{0.01} for any country), an erratum is published identifying the
    affected countries and scores.  The original snapshot and its integrity
    chain remain unchanged.
\end{itemize}

\subsection{Publication Schedule}\label{sec:gov:schedule}

Snapshots are published annually.  The target publication window is
Q2 of the calendar year, reflecting the typical lag in source-data
availability.  The precise date is announced in the project repository at
least 30~days in advance.

\subsection{Archival}\label{sec:gov:archival}

All published snapshots, methodology documents, changelogs, public keys, and
integrity records are archived in a public repository with persistent
identifiers.  No published artefact may be deleted or overwritten.
Corrections are published as addenda, not as replacements.